\section{Soundscape intervention protocol}
\label{sec:protocol}

Figure~\ref{fig:conditions} previews the response profile of the interventions.
Each soundscape contains a localized query event, a temporally separate
nuisance event, and continuous ambience. One of the two answer options names
the query class. The other names a class absent from the placed foregrounds,
and neither option names the nuisance event.

The different-class variant replaces the query event with a recording from the
alternative class. Its matched same-class control inserts a different recording
of the query class into exactly the same interval. Both variants use the same
donor-fitting, level-estimation, background-copying, and crossfade procedures.
Their comparison measures the response to changing the queried class relative
to ordinary rewriting of that interval.

Donor segments remain within their annotated event intervals, and separate
foreground events are at least 0.3\,s apart. The same-class and different-class
replacements use the same target interval and matched edit durations. Post-render
checks verify a median background-noise-floor difference below 1\,dB between
replacement and original spans, fewer than 5\% of soundscapes above a 10\,dB
difference, and boundary first differences within three times those of the
original.

The original audio anchors the six interventions. The sham self-replacement
follows the same editing path, with floating-point rounding changing at most
1 LSB in about 6 of 160,000 samples. The same-class rewrite preserves the
queried class, and whole-clip silence removes the soundscape. The different-class
rewrite substitutes an event from the alternative class. Nuisance removal and
an event-free background-window rewrite test edits away from the query.
Each waveform is queried in both option orders with
``Which of these two sounds is present in it?'', requesting
\texttt{Answer: A} or \texttt{Answer: B}.
